\begin{center}
    {\large \textbf{Giorno 9} \textbar\ Il giorno dell'indipendenza \textbar\ 17 Agosto 2024 \textbar\ \textbf{Gili Trawangan}, Lombok}
\end{center}

\noindent\rule{\textwidth}{0.4pt}
\begin{figure}[h!] % Portrait images below
    \centering
    % First portrait image (on the left)
    \begin{minipage}[h]{0.48\textwidth} % Set width equal to half the page
        \centering
        \includegraphics[width=\textwidth]{images/flags/Screenshot GiliT.png}
    \end{minipage}
    \hfill
    % Text
    \begin{minipage}[h]{0.48\textwidth}
        \raggedright
        \textbf{Superficie}: 4.567 km² \\
        \vspace{0.3cm}
        \textbf{Popolazione}: 3,3 milioni \\
        \vspace{0.3cm}
        \textit{L'isola}
    \end{minipage}
\end{figure}

\landscapeLargeMargin{images/day9/PXL_20240817_005218715.MP.jpg}

\landscapeLargeMargin{images/day9/PXL_20240817_005322379.jpg}

\landscapeLargeMargin{images/day9/PXL_20240817_014441371.jpg}

\landscapeLargeMargin{images/day9/PXL_20240817_020754584.MP.jpg}

\landscapeLargeMargin{images/day9/PXL_20240817_022252428.jpg}

\landscapeLargeMargin{images/day9/PXL_20240817_022624386.MP.jpg}

\landscapeLargeMargin{images/day9/PXL_20240817_085141809.jpg}

\noindent 17 Agosto 2024 \newline

Anche in questa villetta si riposa benissimo! Letto comodo, stanza buia e ben isolata dalla moschea — o comunque dai ripetitori che riproducono preghiere anche al mattino presto — dormiamo tutta una notte filata.

Al risveglio la prima cosa che faccio è ordinare la colazione, o meglio, le colazioni. Elisa la sera prima, mentre io ormai dormivo, ha consultato tutto il menù riga per riga e mi ha scritto un messaggio con la sua ordinazione. Io scelgo i miei due piatti al mattino con gli occhi ancora accartocciati. Tempo una quindicina di minuti e ricevo un messaggio di risposta: il mio ordine è stato preso in carico e sta per arrivare. La gioia.

Ci alziamo, ci vestiamo e prepariamo gli zaini per la giornata: il piano è andare sulle spiagge nel lato ovest dell'isola, che dovrebbero essere le più tranquille poiché situate dalla parte opposta rispetto al centro festaiolo sulla costa est.

Quando arriva la colazione siamo pronti a godercela tutta: due piatti di colazione francese con pain au chocolat, pain au raisin, brioche e pane con marmellata di fragole, un wrap salato per Elisa e un piatto "Brekkie" per me, con uova strapazzate, funghi, pomodori, spinaci e pane tostato. Il tutto accompagnato da caffè di Lombok e succo d'arancia. Io divoro le mie due colazioni dall'inizio alla fine, Eli conserva invece la colazione francese per la merenda del pomeriggio. Trovo assurdo che il caffè di Lombok sia compreso nella colazione mentre per il caffè americano occorra pagare un extra: se fosse stato al contrario avremmo probabilmente pagato l'extra per avere il Lombok coffee.

Terminata la colazione avvisiamo tramite WhatsApp il personale dell'albergo di venire a ritirare i piatti — e di farci il pieno di birra — e usciamo per incontrare Enri e Cri.

Durante il tragitto incontro il mio nuovo migliore amico sull'isola: un gatto arancione, che chiameremo Elisabetta, con una passione irresistibile per le tibie. Molto silenziosamente si posiziona davanti alle mie tibie non appena mi fermo, e per questo motivo si prende sempre un sacco di calci.

Riuniti tutti e quattro ci dirigiamo verso la spiaggia e ci imbattiamo nei primi festeggiamenti: oggi è il 17 agosto, giorno dell'indipendenza indonesiana. Lasciamo i locali a festeggiare tra loro e camminiamo lungo la costa in cerca di un posto dove fermarci. Non ci serve andare molto lontano. In un tratto di spiaggia libera, nascosta sotto gli alberi, troviamo un anfratto spazioso dove coricarci — chi all'ombra, chi al sole — posare gli zaini e stare tranquilli, lontani dal caos e accompagnati solo dal rumore delle onde. Scattiamo qualche foto al mare azzurro davanti a noi e alla spiaggia, approfittando anche di un'altalena appesa a un albero molto scenografica.

Io ed Elisa siamo i primi a mettere la maschera per fare snorkeling. Lei insiste perché io metta le scarpette da roccia, ma non mi lascio certo dire come vestirmi da mia moglie! Dopo aver camminato dieci metri torno indietro sconfitto perché rischio seriamente di tagliarmi i piedi, e mi metto le scarpette da scoglio.

Anche qui i fondali sono spettacolari, ricordano quelli della Coral Beach su Gili Asahan ma più selvaggi e molto più vasti. Siamo attrezzati con action cam e macchina fotografica digitale acquatica e iniziamo a fare foto e video a tutto spiano. A un certo punto Elisa mi chiede di filmarla mentre si immerge. Mi metto di fianco a lei mentre prende fiato e si prepara al tuffo, ma la blocco improvvisamente: dietro di lei sta nuotando una tartaruga gigante! Eli è un po' infastidita per essere stata interrotta proprio mentre aveva appena messo la testa in acqua, ma appena le indico la tartaruga impazzisce di gioia insieme a me. Ci mettiamo a inseguirla scattando foto e facendo video a più non posso. Incrociamo anche un paio di sub che cercano di starle dietro con le pinne, ma si tratta di turisti inesperti e la tartaruga sfugge via a tutti noi. Peccato: non siamo stati molto con lei, ma abbiamo abbastanza materiale fotografico a testimoniare l'incontro.

Torniamo soddisfatti e raccontiamo dell'incontro a Enrico e Cristina, che seguono a ruota il nostro esempio — anche se purtroppo non saranno così fortunati.

Nel frattempo il mare si sta ritirando e i fondali di roccia e coralli vengono allo scoperto, rendendo il paesaggio un po' meno paradisiaco ma permettendo a Eli di andare in esplorazione. Trascorriamo qualche ora tranquilla: chi passeggia per i dintorni, chi girovaga sulla spiaggia alla ricerca di animali e conchiglie, chi sonnecchia. Io sono talmente in pace che riesco quasi ad addormentarmi all'ombra dell'albero, cullato dal suono delle onde. Per pranzo ci teniamo molto leggeri — sei involtini primavera in quattro — e presto giunge il momento di andare a procacciarsi del cibo vero.

Abbandoniamo la nostra caletta e puntiamo verso il sud dell'isola, verso una zona ancora inesplorata. Ovviamo alla fame con un gelato e qualche snack da conservare per l'aperitivo, e nel frattempo ci addentriamo in un'area molto caotica: musica da discoteca, feste sulla spiaggia e addirittura uno schiuma party in un resort con piscina a bordo strada, da cui arrivano nubi di schiuma che decorano la strada in modo molto coreografico. C'è anche un trampolino alto dieci metri dal quale si tuffano un sacco di ragazzini scatenati.

Camminiamo ancora un po', abbastanza decisi a fare dietrofront per tornare nella quiete della villa, quando veniamo attratti da un gruppo di locali che sta facendo una gara di corsa coi sacchi sulla spiaggia. Ci avviciniamo per guardare e scattare qualche foto e immediatamente ci invitano a partecipare. Enri e Cri sono titubanti, ma io ed Elisa non ce lo facciamo ripetere due volte. In realtà speravamo entrambi che ce lo chiedessero.

Ci troviamo nel bel mezzo di una festa locale per il giorno dell'indipendenza indonesiana e scopriamo che hanno organizzato una serie di giochi a premi che ripercorrono i classici giochi della tradizione. Partecipo per primo alla corsa coi sacchi, che si rivela difficilissima: con le mani devo tenere il sacco perché non cada e questo non mi permette di muovere le braccia per bilanciarmi durante i balzi, così devo procedere a saltelli e rischio più volte di cadere rovinosamente, per fortuna senza perdere troppo la dignità. Dopo è il turno di Elisa che, dopo aver visto la tartaruga, ne ha fatto il suo animale guida. Diciamo che arriva a termine corsa.

Prima di iniziare la sfida successiva ci invitano a banchettare con prodotti locali: crocchette di verdure preparate con la pasta dell'involtino primavera, da intingere nella salsa piccante, e dolci di riso avvolti nella foglia di banano. Da bere un tè dolce allo zenzero preparato da loro, di cui temevamo l'origine del ghiaccio ma che non disdegnamo.

Il secondo round di giochi è il tiro alla fune e convinciamo anche Enrico e Cristina a unirsi. La prima sfida è tra la mia squadra e quella di Enrico. Il secondo membro della mia squadra è più basso di me di almeno trenta centimetri, mentre la squadra di Enrico è messa un po' meglio. Chiaramente queste sono le facili scuse a cui cerco di appigliarmi, dal momento che ci asfaltano velocemente. D'altronde Enrico è figlio d'arte: suo padre era campione italiano di tiro alla fune e lui si è sempre cimentato alle feste di paese con suo fratello fin da piccoli. Altre scuse. La sua squadra passa il turno ma viene poi battuta in semifinale.

Terminata la gara degli uomini è toccato alle donne: solo due squadre stavolta, con Elisa e Cristina l'una contro l'altra. Per Elisa è stata la rivalsa: hanno sconfitto le avversarie in entrambi i round.

Questi ragazzi sono stati davvero aperti e calorosi a invitarci a festeggiare con loro nel loro giorno dell'indipendenza, a condividere il loro cibo e a farci conoscere le loro tradizioni. Li ringraziamo di cuore, ma si sta facendo tardi e siamo costretti a salutarli. Tornati nella villa stappiamo qualche birra, condividiamo gli snack acquistati nel pomeriggio e ci rilassiamo un'oretta, ricaricando le energie. E le macchine fotografiche.

Terminato l'aperitivo i nostri amici tornano al loro hotel mentre noi ci facciamo la doccia e ci prepariamo per uscire. Durante i preparativi Elisa scorge nel bagno, sulla terra dove crescono i bambù, una chiocciola gigantesca, lunga quasi quanto la mia mano aperta.

Docciati e vestiti, ci incamminiamo alla ricerca della cena. Attraversando le strade buie di Trawangan ci arriva spesso un odore acre di bruciato: purtroppo non si tratta solo di fogliame, ma anche di gomma, il cui fumo denso irrita gola e narici. Ogni sera, nei giardini delle case o nei campi, i rifiuti vengono dati alle fiamme in grandi falò.

Nelle aree interne dell'isola la luce è completamente assente e le famiglie indonesiane cenano quasi tutte al buio. Questo contrasta nettamente con la luminosità e il fermento della via principale lungo la spiaggia: hotel con piscine scintillanti, ristoranti e piano-bar con musica dal vivo, noleggiatori di tavole da surf e attrezzatura da sub, bancarelle di souvenir e tour operator. Il crepitio dei falò, il muggito delle vacche che vagano libere e il richiamo alla preghiera si mescolano al sottofondo di musica commerciale, house tribale e jazz che anima le strade in festa.

Arrivati nel centro della movida ci accorgiamo che cenare non sarà così facile: in molti locali non ci fanno sedere perché hanno già chiuso la cucina. "Just for drinks" è la frase che sentiamo più spesso. Dopo vari tentativi riusciamo finalmente a sederci all'Ocean 2. Dopo il successo della cena precedente, questa sera l'esperienza è tutt'altro che soddisfacente. Ordiniamo involtini primavera, Mie Goreng ai frutti di mare e Curry Soup. Gli involtini sono intrisi d'olio e sul tavolo ci sono solo piccoli tovagliolini trasparenti, del tutto insufficienti a evitare che le mani restino unte. Il mio Mie Goreng arriva con il pollo al posto dei frutti di mare, mentre Elisa lascia da parte tutto il pesce del suo piatto: è andato a male.

L'idea per la serata è di girovagare in cerca di intrattenimento, ma i posti non ci ispirano — o forse abbiamo tutti e quattro gusti un po' diversi. Bazzichiamo davanti alle discoteche senza fermarci e ci sediamo poi in un bar a bere qualcosa. Anche l'idea degli shots fallisce miseramente, poiché la sensazione è che cerchino di propinarci gli alcolici più cari. Per fortuna la gentilezza indonesiana non viene meno neanche stasera e alle mie rimostranze ci cambiano lo scontrino e ci permettono di modificare l'ordinazione. Il mio cocktail — un mojito estremamente zuccherato, alternativa al costosissimo shot — mi restituisce un po' di energie. È così che decidiamo di tornare in uno dei locali con DJ set e proviamo a ballare un po', cercando di evitare gli spacciatori.

Questo è un problema non da poco: pare che sull'isola crescano funghi allucinogeni, che sono il prodotto più venduto insieme alla marijuana. Se non il più venduto, è quello che ci propongono più spesso, soprattutto i tanti venditori ambulanti che probabilmente sono anche piccoli produttori. Per carità, fate pure — chi sono io per giudicare — ma se evitate di chiedermelo otto volte in due ore sono più contento.

Si è fatto ormai tardi e dal momento che l'indomani ci aspetta un'altra giornata di sole decidiamo di tornare alla villa. Durante il tragitto ritrovo il gatto Elisabetta, ma stavolta evito di darle calci.

\clearpage

\textit{P.s. Un estratto del programma dei giochi organizzati in occasione della festa dell'indipendenza indonesiana, che mi hanno gentilmente lasciato fotografare.}
\newline

\textbf{SPIRIT OF INDONESIAN RACE}
\newline

Buon pomeriggio a tutti! 
\newline

Benvenuti alla nostra celebrazione del Giorno dell'Indipendenza, presentata con orgoglio da Lumi Hotel. Oggi abbiamo in programma per voi una serie di entusiasmanti giochi tradizionali. Ci divertiremo con la Corsa alle biglie, la Gara di mangiatori di cracker, la Corsa con i sacchi, il Tiro alla fune, il Chiodo nella bottiglia e la Sfida del trucco. Questo evento è aperto sia ai nostri ospiti che al personale, quindi uniamoci tutti al divertimento e rendiamo questa giornata indimenticabile!
\newline

Inizio evento:
\newline

1. Il nostro primo gioco è la "Corsa delle biglie". In questa competizione i partecipanti utilizzeranno un cucchiaio tenuto in bocca per portare una biglia fino al traguardo. Ci saranno gare separate per uomini e donne e i vincitori accederanno alla fase finale. Vediamo chi ha il miglior equilibrio!
\newline

2. Che gara emozionante è stata! Congratulazioni a tutti i partecipanti e un grande applauso ai nostri vincitori! Ora prepariamoci per la nostra prossima sfida. Sei pronto per divertirti di più?
\newline

3. Il prossimo è il "Concorso di mangiatori di cracker". L'obiettivo è finire di mangiare il cracker il più velocemente possibile senza usare le mani! Preparati a tifare per i tuoi partecipanti preferiti! Wow, che finale al cardiopalmo! Complimenti a tutti coloro che hanno partecipato. Manteniamo vivo questo entusiasmo mentre passiamo alla prossima sfida!
\newline

4. Passiamo ora alla "corsa ai sacchi". I partecipanti salteranno fino al traguardo in un sacco. È una gara di agilità e velocità e abbiamo round separati per uomini e donne. Buona fortuna a tutti! Uno sforzo fantastico da parte di tutti i nostri partecipanti! La competizione si sta infiammando e il meglio deve ancora venire. Sei pronto per il prossimo round?
\newline

5. Il nostro prossimo evento è il "Tiro alla fune". Squadre di cinque persone si sfideranno tirando la corda finché una squadra non taglierà il traguardo. La forza e il lavoro di squadra sono fondamentali, quindi dai il massimo! Incredibile! È stata una gara emozionante. Un applauso per i nostri vincitori e vediamo chi brillerà nella prossima partita!
\newline

6. Nel nostro gioco "Chiodo nella bottiglia", le squadre dovranno lavorare insieme per guidare un chiodo in una bottiglia mentre sono bendati. Solo un membro del team può dare istruzioni, quindi la comunicazione è fondamentale. Vediamo chi può lavorare meglio insieme! Lo spirito di competizione è davvero vivo oggi! Congratulazioni ai nostri vincitori. Vediamo se il prossimo gioco riuscirà a superare quell'entusiasmo!


\pagestyle{empty} % disable numbers

\twoImagesLargeMargins{images/day9/PXL_20240817_005024318.jpg}
                      {images/day9/PXL_20240817_005206599.jpg}
%coppia
\splitImageLeft{images/day9/PXL_20240817_014451812.jpg}
\splitImageRight{images/day9/PXL_20240817_014451812.jpg}
%

%coppia
\splitImageLeft{images/day9/PXL_20240817_020749328.MP.jpg}
\splitImageRight{images/day9/PXL_20240817_020749328.MP.jpg}
%

%coppia
\twoImagesLargeMargins{images/day9/PXL_20240817_023200751.jpg}
                      {images/day9/PXL_20240817_022624386.MP.jpg}
\twoImagesLargeMargins{images/day9/PXL_20240817_023657939.MP.jpg}
                      {images/day9/AAAF3F3E-9B2B-4565-AE5A-5BF4EAC1979F_1_201_a.jpeg}
%

%coppia
\portraitFullscreen{images/day9/PXL_20240817_023600386.jpg}
\landscapeLargeMargin{images/day9/PXL_20240817_023316611.MP.jpg}
%

%coppia
\splitImageLeft{images/day9/PXL_20240817_022316297.MP.jpg}
\splitImageRight{images/day9/PXL_20240817_022316297.MP.jpg}
%

%coppia
\twoImagesLargeMargins{images/day9/PXL_20240817_085335659.jpg}
                      {images/day9/PXL_20240817_093427724.jpg}
\twoImagesLargeMargins{images/day9/PXL_20240817_085932991.jpg}
                      {images/day9/PXL_20240817_085859975.MP.jpg}

%coppia
\sixLandscapeImages{images/day9/PXL_20240817_091728186.MP.jpg}
                   {images/day9/PXL_20240817_091729486.MP.jpg}
                   {images/day9/PXL_20240817_091730336.jpg}
                   {images/day9/PXL_20240817_091734602.MP.jpg}
                   {images/day9/7215386C-2942-419A-A4A7-8EF3EC21257A_1_201_a.jpeg}
                   {images/day9/0AC4D2DD-7442-4507-B2E0-2575C8F7659D_1_201_a.jpeg}

%coppia
\sixLandscapeImages{images/day9/PXL_20240817_092825289.jpg}
                   {images/day9/PXL_20240817_092827839.jpg}
                   {images/day9/73DF18A3-42E5-41E4-850B-914EBBC24590_1_201_a.jpeg}
                   {images/day9/20932FD7-8154-4031-9F8A-E68E0BBAC611_1_201_a.jpeg}
                   {images/day9/253E6D0E-6B75-48BA-8937-7E365FC09F99_1_201_a.jpeg}
                   {images/day9/CC39CE3E-D5B5-4E6C-A8F1-BC9B73601970_1_201_a.jpeg}

%coppia
\twoImagesLargeMargins{images/day9/PXL_20240817_090345510.jpg}
                      {images/day9/PXL_20240817_090351114.jpg}
\twoImagesLargeMargins{images/day9/PXL_20240817_090208521.MP.jpg}
                      {images/day9/PXL_20240817_090145155.jpg} 

%coppia
\twoImagesLargeMargins{images/day9/PXL_20240817_120021703.jpg}
                      {images/day9/PXL_20240817_101202294.MP.jpg}
\portraitFullscreen{images/day9/PXL_20240817_142553850.jpg}
%

\twoImagesLargeMargins{images/day9/PXL_20240817_161912210.jpg}
                      {images/day9/PXL_20240817_161915731.jpg}

\pagestyle{fancy} % enable numbers